\documentclass[11pt,a4paper]{article}
\usepackage{templates/modeling-notes}
\hypersetup{pdftitle={Sheet 1 -- Dual simplex and primal recovery}}
\title{\textbf{Mathematical Modeling}\\[2mm]
\large Sheet 1: Dual Simplex and Primal Recovery\\[1mm]
\normalsize Tasks 4 and 5(a)--5(d)}
\author{Personal study notes}
\date{}
\begin{document}
\maketitle
\tableofcontents

\section{Task statements}\label{sec:statements}
\begin{otherlanguage}{polish}
  \noindent\textbf{4.} Rozwiąż algorytmem dualnym. Na tej podstawie podaj rozwiązanie zagadnienie dualnego.
  \[
    \begin{aligned}
      \min\quad &x_1+6x_2+2x_3+3x_4,\\
      &4x_1+2x_2-x_3+x_4\geq4,\\
      &x_1+x_2+x_3+x_4\geq5,\\
      &-2x_1-2x_2-x_3-x_4\geq-8,\\
      &x_1,x_2,x_3,x_4\geq0.
    \end{aligned}
  \]
  \noindent\textbf{5.}
  \begin{itemize}[leftmargin=*]
    \item Podaj postać zadania dualnego ZD
    \item Rozwiąż ZD dwufazową metodą sympleks i podaj jego znalezione RO oraz optymalną wartość funkcji celu dla ZD
    \item Na podstawie znalezionego rozwiązania RO dla zadania ZD, oblicz RO dla zadania ZP ( bez rozwiązywania ZP) oraz sprawdź optymalną wartość funkcji w ZP
    \item Czy zadanie ZD posiada rozwiązania alternatywne? Uzasadnij
  \end{itemize}
  \noindent\textbf{(a)}
  \[
    \begin{aligned}
      \min\quad &3x_1+10x_2+2x_3,\\
      &-x_1+2x_2+x_3\geq2,\\
      &-x_1+x_2+x_3\geq1,\\
      &-x_1+2x_2-x_3\geq-3,\\
      &-x_1+2x_2+x_3\geq1,\\
      &x_1\leq0,\quad x_2\geq0,\quad x_3\geq0.
    \end{aligned}
  \]
  \noindent\textbf{(b)}
  \[
    \begin{aligned}
      \min\quad &3x_1+10x_2+x_3,\\
      &-2x_2-2x_3\geq-2,\\
      &2x_1-x_2\geq1,\\
      &-2x_1+x_2+x_3\geq-5,\\
      &2x_1+2x_2+x_3\geq1,\\
      &x_1\geq0,\quad x_2\in\R,\quad x_3\geq0.
    \end{aligned}
  \]
  \noindent\textbf{(c)}
  \[
    \begin{aligned}
      \min\quad &-x_1+5x_2+8x_3,\\
      &-x_1-x_2+2x_3\geq-1,\\
      &-x_1+3x_3\geq5,\\
      &-x_1+x_2+x_3\geq1,\\
      &-x_2-2x_3\geq-1,\\
      &x_1\in\R,\quad x_2\in\R,\quad x_3\geq0.
    \end{aligned}
  \]
  \noindent\textbf{(d)}
  \[
    \begin{aligned}
      \min\quad &x_1-5x_2+8x_3,\\
      &3x_1+x_2+2x_3\geq1,\\
      &-x_1+x_3\geq5,\\
      &-x_1-x_2+x_3\geq1,\\
      &3x_1-x_2-2x_3\geq-2,\\
      &x_1\geq0,\quad x_2\in\R,\quad x_3\geq0.
    \end{aligned}
  \]
\end{otherlanguage}

\section{Notation, domains, and the mathematical plan}\label{sec:conventions}
All decision variables are continuous and obey the displayed sign restrictions.
The data specify abstract linear programs, with no physical units supplied.
In particular, a free variable may be negative: it must not acquire an extra
nonnegativity restriction during conversion or verification.

The prerequisites are standard-form equivalence and bases in
\href{../lectures/02-linear-programming.pdf}{Chapter 02}, Phase I and reduced costs in
\href{../lectures/03-simplex-method.pdf}{Chapter 03}, and mixed-domain duality,
complementarity, and dual simplex in
\href{../lectures/04-degeneracy-and-duality.pdf}{Chapter 04}.
We apply these ideas to the complete tasks below.

For any equality-form maximization, use
\[
  \max c^Tu,\qquad Mu=d,\quad u\geq0.
\]
A nonsingular basis $B$ gives
\[
  \bar M=M_B^{-1}M,\quad \bar d=M_B^{-1}d,\quad
  \lambda=M_B^{-T}c_B,\quad q=M^T\lambda-c.
\]
The tableau displays the equations $\bar M u=\bar d$ and the last row
$q^Tu=z_0-c^Tu$, where $z_0=c_B^T\bar d$. Thus
\begin{equation}\label{eq:indicator}
  c^Tu=z_0-\sum_{j\notin B}q_ju_j.
\end{equation}
The last RHS is $z_0$, not a basic variable. Ratios are calculated in the column to the right of the RHS;
$\longrightarrow$ in the separate Out column marks the chosen leaving row.
The indicator row is labelled $z_j-c_j$; $q_j$ abbreviates this difference
in formulas. Vertical rules separate $c_B$, the basic-variable labels,
the coefficient columns, and the RHS. Variable-column headings keep their
order across pivots; only the basic-variable labels change. Artificial
columns are removed after Phase I without reordering the retained columns. The $c_B$ column gives the objective coefficient of each basic
variable. The $z_j$ row sums each constraint column with these weights;
subtract the displayed $c_j$ row to obtain $q_j=z_j-c_j$. The same
weighted sum on the RHS gives $z_0$. In Phase I these are the auxiliary
objective coefficients; in Phase II they are the original ones. Nonnegative indicators certify
an upper bound on the maximum. The original minimum in task 4 is the negative
of its transformed maximum. Tasks 5 maximize their dual objectives directly.

\subsection{Why the domains determine the dual inequalities}
Write a primal minimum as $\min c^Tx$, $Ax\geq b$. A nonnegative multiplier
$y$ gives $y^TAx\geq y^Tb$. To make $c^Tx\geq y^TAx$ hold for every allowed $x$,
column $j$ must obey
\[
  \begin{array}{c|c}
    \text{Primal domain}&\text{Dual condition}\\\hline
    x_j\geq0&a_j^Ty\leq c_j\\
    x_j\leq0&a_j^Ty\geq c_j\\
    x_j\in\R&a_j^Ty=c_j
  \end{array}
\]
For example, multiplying a nonpositive $x_j$ by $c_j-a_j^Ty\leq0$ gives a
nonnegative product. A free variable can take either sign, so its coefficient
must vanish. The resulting dual maximizes $b^Ty$ over $y\geq0$ with these
column conditions. Negative entries in $b$ do not change multiplier domains.

For a feasible pair, the gap has the decomposition
\begin{equation}\label{eq:gap}
  c^Tx-b^Ty=(c-A^Ty)^Tx+y^T(Ax-b)\geq0.
\end{equation}
Every term on the right is nonnegative with the appropriate domain convention.
Optimal pairs make each term zero. This explains complementary slackness,
rather than treating it as an additional system to memorize.

\subsection{Two different simplex invariants}
Primal simplex keeps $\bar d\geq0$. It chooses an entering column with $q_k<0$
and preserves feasibility by the ratio test
\[
  \theta=\min_{\bar m_{ik}>0}\frac{\bar d_i}{\bar m_{ik}}.
\]
Our Phase I and II paths choose the most negative indicator, breaking ties by
smaller column index, then the smallest ratio, breaking ties by smaller leaving
variable index. These choices specify these finite paths; they are not a general
anti-cycling guarantee.

Dual simplex keeps $q\geq0$. It first selects a negative RHS row $r$.
Eligible entering columns have $\bar m_{rk}<0$, and
\[
  \alpha=\min_{\bar m_{rj}<0}\frac{q_j}{-\bar m_{rj}}
\]
preserves nonnegative indicators because $q'_j=q_j+\alpha\bar m_{rj}$.
We select the most negative RHS and the smaller entering column in a tie.
The update $z'_0=z_0+\alpha\bar d_r$ lowers the dual upper bound on a maximum.
It does not describe improvement of an infeasible primal point.

Phase I maximizes minus the sum of nonnegative artificial variables.
Its optimum is at most zero. A negative optimum certifies infeasibility; zero
permits Phase II after artificials are removed and the original objective is
recomputed. All successful paths here remove artificials from the basis, so no
zero-basic artificial cleanup is needed.

\section{Task 4: a complete dual-simplex solve}\label{sec:task4}
Put $s_i=(Ax-b)_i\geq0$, negate each row, and maximize the negative cost:
\[
  \begin{aligned}
    \max\quad &-x_1-6x_2-2x_3-3x_4,\\
    -4x_1-2x_2+x_3-x_4+s_1&=-4,\\
    -x_1-x_2-x_3-x_4+s_2&=-5,\\
    2x_1+2x_2+x_3+x_4+s_3&=8.
  \end{aligned}
\]
The surplus basis $B=(s_1,s_2,s_3)$ is nonsingular, with indicators
$(1,6,2,3,0,0,0)\geq0$. Its negative RHS entries make it primal infeasible,
but exactly the dual-feasible start we need.

\begin{center}
  \small
  \renewcommand{\arraystretch}{1.25}
  \begin{tabular}{r@{\hspace{0.2em}}|@{\hspace{0.2em}}c@{\enspace}|@{\enspace}rrrrrrr@{\quad}|@{\quad}r@{\qquad\quad}l@{\quad}c}\cmidrule(r{1.5em}){1-10}
    $c_B$ & Basic & $x_1$ & $x_2$ & $x_3$ & $x_4$ & $s_1$ & $s_2$ & $s_3$ & RHS & Dual ratios ($r=2$) & Out\\
    & $c_j$ & $-1$ & $-6$ & $-2$ & $-3$ & $0$ & $0$ & $0$ & & & \\
    \cmidrule(r{1.5em}){1-10}
    $0$ & $s_1$ & $-4$ & $-2$ & $1$ & $-1$ & $1$ & $0$ & $0$ & $-4$ & $x_1:\ 1/(1)=1$ & \\
    $0$ & $s_2$ & $-1$ & $-1$ & $-1$ & $-1$ & $0$ & $1$ & $0$ & $-5$ & $x_2:\ 6/(1)=6$ & $\longrightarrow$\\
    $0$ & $s_3$ & $2$ & $2$ & $1$ & $1$ & $0$ & $0$ & $1$ & $8$ & $x_3:\ 2/(1)=2,\quad x_4:\ 3/(1)=3$ & \\
    \cmidrule(r{1.5em}){1-10}
    & $z_j$ & $0$ & $0$ & $0$ & $0$ & $0$ & $0$ & $0$ & $0$ & & \\
    & $z_j-c_j$ & $1$ & $6$ & $2$ & $3$ & $0$ & $0$ & $0$ & $0$ & & \\
    \cmidrule(r{1.5em}){1-10}
  \end{tabular}
\end{center}

The most negative RHS is $-5$: $s_2$ leaves and $x_1$ enters on the pivot $-1$.
Divide row 2 by $-1$, then eliminate column 1 from rows 1 and 3. The new bound
is $z_0=0+1(-5)=-5$.

\begin{center}
  \small
  \renewcommand{\arraystretch}{1.25}
  \begin{tabular}{r@{\hspace{0.2em}}|@{\hspace{0.2em}}c@{\enspace}|@{\enspace}rrrrrrr@{\quad}|@{\quad}r@{\qquad\quad}l@{\quad}c}\cmidrule(r{1.5em}){1-10}
    $c_B$ & Basic & $x_1$ & $x_2$ & $x_3$ & $x_4$ & $s_1$ & $s_2$ & $s_3$ & RHS & Dual ratios ($r=3$) & Out\\
    & $c_j$ & $-1$ & $-6$ & $-2$ & $-3$ & $0$ & $0$ & $0$ & & & \\
    \cmidrule(r{1.5em}){1-10}
    $0$ & $s_1$ & $0$ & $2$ & $5$ & $3$ & $1$ & $-4$ & $0$ & $16$ & $x_3:\ 1/(1)=1$ & \\
    $-1$ & $x_1$ & $1$ & $1$ & $1$ & $1$ & $0$ & $-1$ & $0$ & $5$ & $x_4:\ 2/(1)=2$ & \\
    $0$ & $s_3$ & $0$ & $0$ & $-1$ & $-1$ & $0$ & $2$ & $1$ & $-2$ & & $\longrightarrow$\\
    \cmidrule(r{1.5em}){1-10}
    & $z_j$ & $-1$ & $-1$ & $-1$ & $-1$ & $0$ & $1$ & $0$ & $-5$ & & \\
    & $z_j-c_j$ & $0$ & $5$ & $1$ & $2$ & $0$ & $1$ & $0$ & $-5$ & & \\
    \cmidrule(r{1.5em}){1-10}
  \end{tabular}
\end{center}

Row 3 still has RHS $-2$, so $s_3$ leaves and $x_3$ enters on the pivot $-1$.
The next bound is $-5+1(-2)=-7$.

\begin{center}
  \small
  \renewcommand{\arraystretch}{1.25}
  \begin{tabular}{r@{\hspace{0.2em}}|@{\hspace{0.2em}}c@{\enspace}|@{\enspace}rrrrrrr@{\quad}|@{\quad}r}\toprule
    $c_B$ & Basic & $x_1$ & $x_2$ & $x_3$ & $x_4$ & $s_1$ & $s_2$ & $s_3$ & RHS\\
    & $c_j$ & $-1$ & $-6$ & $-2$ & $-3$ & $0$ & $0$ & $0$ & \\
    \midrule
    $0$ & $s_1$ & $0$ & $2$ & $0$ & $-2$ & $1$ & $6$ & $5$ & $6$\\
    $-1$ & $x_1$ & $1$ & $1$ & $0$ & $0$ & $0$ & $1$ & $1$ & $3$\\
    $-2$ & $x_3$ & $0$ & $0$ & $1$ & $1$ & $0$ & $-2$ & $-1$ & $2$\\
    \midrule
    & $z_j$ & $-1$ & $-1$ & $-2$ & $-2$ & $0$ & $3$ & $1$ & $-7$\\
    & $z_j-c_j$ & $0$ & $5$ & $0$ & $1$ & $0$ & $3$ & $1$ & $-7$\\
    \bottomrule
  \end{tabular}
\end{center}

Now both invariants hold: every RHS and indicator is nonnegative. Reading
nonbasic variables as zero gives
\[
  x^*=(3,0,2,0),\quad s^*=(6,0,0),\qquad \min c^Tx=7.
\]
The final tableau also gives the exact identity
\begin{equation}\label{eq:task4-bound}
  c^Tx=7+5x_2+x_4+3s_2+s_3\geq7.
\end{equation}
Equality forces $x_2=x_4=s_2=s_3=0$. The last two tight primal rows then give
$x_1+x_3=5$ and $2x_1+x_3=8$, so this primal optimizer is unique.

\subsection{Read and verify the dual certificate}
The original minimum has dual
\[
  \begin{aligned}
    \max\quad &4y_1+5y_2-8y_3,\\
    4y_1+y_2-2y_3&\leq1,\\
    2y_1+y_2-2y_3&\leq6,\\
    -y_1+y_2-y_3&\leq2,\\
    y_1+y_2-y_3&\leq3,\qquad y\geq0.
  \end{aligned}
\]
For the transformed equations $[-A,I](x,s)=-b$, the basis multiplier satisfies
$-A^Ty\geq-c$ and $y\geq0$, which are precisely these constraints.
The last tableau's original surplus indicators give
\[
  y^*=(0,3,1),\quad c-A^Ty^*=(0,5,0,1),\quad b^Ty^*=15-8=7.
\]
We have primal and dual feasibility, equal values, and
$y^*\odot s^*=0$, $x^*\odot(c-A^Ty^*)=0$.
To see dual uniqueness, $s_1=6>0$ forces $y_1=0$ in every dual optimum.
The positive $x_1$ and $x_3$ force $y_2-2y_3=1$ and $y_2-y_3=2$,
whose unique solution is $(y_2,y_3)=(3,1)$.

\section{Task 5(a): Phase I proves dual infeasibility}\label{sec:task5a}
The nonpositive domain of $x_1$ reverses its dual inequality. The dual is
\[
  \begin{aligned}
    \max\quad &2y_1+y_2-3y_3+y_4,\\
    -y_1-y_2-y_3-y_4&\geq3,\\
    2y_1+y_2+2y_3+2y_4&\leq10,\\
    y_1+y_2-y_3+y_4&\leq2,\qquad y\geq0.
  \end{aligned}
\]
Its first condition is impossible because its left side is nonpositive.
We can also express this contradiction through the requested two-phase procedure.
Introduce surplus $t_1$ and artificial $a_1$ in row 1, and slacks $s_2,s_3$:
\[
  -y_1-y_2-y_3-y_4-t_1+a_1=3.
\]
Use $B=(a_1,s_2,s_3)$ and maximize $-a_1$. The other two equations have
right sides $10$ and $2$.

\begin{center}
  \small
  \renewcommand{\arraystretch}{1.25}
  \begin{tabular}{r@{\hspace{0.2em}}|@{\hspace{0.2em}}c@{\enspace}|@{\enspace}rrrrrrrr@{\quad}|@{\quad}r}\toprule
    $c_B$ & Basic & $y_1$ & $y_2$ & $y_3$ & $y_4$ & $t_1$ & $s_2$ & $s_3$ & $a_1$ & RHS\\
    & $c_j$ & $0$ & $0$ & $0$ & $0$ & $0$ & $0$ & $0$ & $-1$ & \\
    \midrule
    $-1$ & $a_1$ & $-1$ & $-1$ & $-1$ & $-1$ & $-1$ & $0$ & $0$ & $1$ & $3$\\
    $0$ & $s_2$ & $2$ & $1$ & $2$ & $2$ & $0$ & $1$ & $0$ & $0$ & $10$\\
    $0$ & $s_3$ & $1$ & $1$ & $-1$ & $1$ & $0$ & $0$ & $1$ & $0$ & $2$\\
    \midrule
    & $z_j$ & $1$ & $1$ & $1$ & $1$ & $1$ & $0$ & $0$ & $-1$ & $-3$\\
    & $z_j-c_j$ & $1$ & $1$ & $1$ & $1$ & $1$ & $0$ & $0$ & $0$ & $-3$\\
    \bottomrule
  \end{tabular}
\end{center}

Every indicator is already nonnegative. Therefore Phase I has optimum $-3$,
or minimum artificial sum $w=3$. More directly,
\[
  a_1=3+y_1+y_2+y_3+y_4+t_1\geq3.
\]
A feasible original dual would need $a_1=0$, which is impossible. Stop before
Phase II: there is no dual optimum and hence no optimal pair to recover.

\subsection{A primal ray explains the same outcome}
The ray
\[
  x(t)=(-2-t,0,0),\qquad t\geq0
\]
respects all domains and has primal row values $(2+t,2+t,2+t,2+t)$,
which dominate $(2,1,-3,1)$. Its objective is $-6-3t\to-\infty$.
Thus the primal is feasible and unbounded below. This is consistent with the
dual contradiction: a feasible dual would supply a finite lower bound and
would forbid such a ray. Neither problem has an optimal solution;
``alternative optimal dual solutions'' therefore do not exist.

\section{Task 5(b): solve the dual and recover the primal}\label{sec:task5b}
Because $x_2$ is free, its column gives an equality:
\[
  \begin{aligned}
    \max\quad &-2y_1+y_2-5y_3+y_4,\\
    2y_2-2y_3+2y_4&\leq3,\\
    -2y_1-y_2+y_3+2y_4&=10,\\
    -2y_1+y_3+y_4&\leq1,\qquad y\geq0.
  \end{aligned}
\]
Let $s_1,s_3$ be the slacks of the two dual inequalities, and add $a_2$ to
its equality. Column order is $(y_1,y_2,y_3,y_4,s_1,s_3,a_2)$.
Notice that these slacks index dual conditions, not primal resource rows.

\subsection{Phase I: three pivots remove the artificial variable}
Maximize $-a_2$, starting at $B=(s_1,a_2,s_3)$.

\begin{center}
  \small
  \renewcommand{\arraystretch}{1.25}
  \begin{tabular}{r@{\hspace{0.2em}}|@{\hspace{0.2em}}c@{\enspace}|@{\enspace}rrrrrrr@{\quad}|@{\quad}r@{\qquad\quad}l@{\quad}c}\cmidrule(r{1.5em}){1-10}
    $c_B$ & Basic & $y_1$ & $y_2$ & $y_3$ & $y_4$ & $s_1$ & $s_3$ & $a_2$ & RHS & Ratio ($y_4\uparrow$) & Out\\
    & $c_j$ & $0$ & $0$ & $0$ & $0$ & $0$ & $0$ & $-1$ & & & \\
    \cmidrule(r{1.5em}){1-10}
    $0$ & $s_1$ & $0$ & $2$ & $-2$ & $2$ & $1$ & $0$ & $0$ & $3$ & $3 / 2$ & \\
    $-1$ & $a_2$ & $-2$ & $-1$ & $1$ & $2$ & $0$ & $0$ & $1$ & $10$ & $10 / 2=5$ & \\
    $0$ & $s_3$ & $-2$ & $0$ & $1$ & $1$ & $0$ & $1$ & $0$ & $1$ & $1 / 1=1$ & $\longrightarrow$\\
    \cmidrule(r{1.5em}){1-10}
    & $z_j$ & $2$ & $1$ & $-1$ & $-2$ & $0$ & $0$ & $-1$ & $-10$ & & \\
    & $z_j-c_j$ & $2$ & $1$ & $-1$ & $-2$ & $0$ & $0$ & $0$ & $-10$ & & \\
    \cmidrule(r{1.5em}){1-10}
  \end{tabular}
\end{center}

The most negative indicator is $q_4=-2$. Enter $y_4$ and remove
$s_3$ after step $1$, raising the
Phase I value from $-10$ to $-8$.

\begin{center}
  \small
  \renewcommand{\arraystretch}{1.25}
  \begin{tabular}{r@{\hspace{0.2em}}|@{\hspace{0.2em}}c@{\enspace}|@{\enspace}rrrrrrr@{\quad}|@{\quad}r@{\qquad\quad}l@{\quad}c}\cmidrule(r{1.5em}){1-10}
    $c_B$ & Basic & $y_1$ & $y_2$ & $y_3$ & $y_4$ & $s_1$ & $s_3$ & $a_2$ & RHS & Ratio ($y_1\uparrow$) & Out\\
    & $c_j$ & $0$ & $0$ & $0$ & $0$ & $0$ & $0$ & $-1$ & & & \\
    \cmidrule(r{1.5em}){1-10}
    $0$ & $s_1$ & $4$ & $2$ & $-4$ & $0$ & $1$ & $-2$ & $0$ & $1$ & $1 / 4$ & $\longrightarrow$\\
    $-1$ & $a_2$ & $2$ & $-1$ & $-1$ & $0$ & $0$ & $-2$ & $1$ & $8$ & $8 / 2=4$ & \\
    $0$ & $y_4$ & $-2$ & $0$ & $1$ & $1$ & $0$ & $1$ & $0$ & $1$ & $\text{--}$ & \\
    \cmidrule(r{1.5em}){1-10}
    & $z_j$ & $-2$ & $1$ & $1$ & $0$ & $0$ & $2$ & $-1$ & $-8$ & & \\
    & $z_j-c_j$ & $-2$ & $1$ & $1$ & $0$ & $0$ & $2$ & $0$ & $-8$ & & \\
    \cmidrule(r{1.5em}){1-10}
  \end{tabular}
\end{center}

Now $y_1$ enters because $q_1=-2$; $s_1$ leaves after step $1/4$.

\begin{center}
  \small
  \renewcommand{\arraystretch}{1.25}
  \begin{tabular}{r@{\hspace{0.2em}}|@{\hspace{0.2em}}c@{\enspace}|@{\enspace}rrrrrrr@{\quad}|@{\quad}r@{\qquad\quad}l@{\quad}c}\cmidrule(r{1.5em}){1-10}
    $c_B$ & Basic & $y_1$ & $y_2$ & $y_3$ & $y_4$ & $s_1$ & $s_3$ & $a_2$ & RHS & Ratio ($y_3\uparrow$) & Out\\
    & $c_j$ & $0$ & $0$ & $0$ & $0$ & $0$ & $0$ & $-1$ & & & \\
    \cmidrule(r{1.5em}){1-10}
    $0$ & $y_1$ & $1$ & $\tfrac{1}{2}$ & $-1$ & $0$ & $\tfrac{1}{4}$ & $-\tfrac{1}{2}$ & $0$ & $\tfrac{1}{4}$ & $\text{--}$ & \\
    $-1$ & $a_2$ & $0$ & $-2$ & $1$ & $0$ & $-\tfrac{1}{2}$ & $-1$ & $1$ & $\tfrac{15}{2}$ & $(15 / 2) / 1=15 / 2$ & $\longrightarrow$\\
    $0$ & $y_4$ & $0$ & $1$ & $-1$ & $1$ & $\tfrac{1}{2}$ & $0$ & $0$ & $\tfrac{3}{2}$ & $\text{--}$ & \\
    \cmidrule(r{1.5em}){1-10}
    & $z_j$ & $0$ & $2$ & $-1$ & $0$ & $\tfrac{1}{2}$ & $1$ & $-1$ & $-\tfrac{15}{2}$ & & \\
    & $z_j-c_j$ & $0$ & $2$ & $-1$ & $0$ & $\tfrac{1}{2}$ & $1$ & $0$ & $-\tfrac{15}{2}$ & & \\
    \cmidrule(r{1.5em}){1-10}
  \end{tabular}
\end{center}

Finally $y_3$ enters, with indicator $-1$, and the artificial $a_2$
leaves after step $15/2$.

\begin{center}
  \small
  \renewcommand{\arraystretch}{1.25}
  \begin{tabular}{r@{\hspace{0.2em}}|@{\hspace{0.2em}}c@{\enspace}|@{\enspace}rrrrrrr@{\quad}|@{\quad}r}\toprule
    $c_B$ & Basic & $y_1$ & $y_2$ & $y_3$ & $y_4$ & $s_1$ & $s_3$ & $a_2$ & RHS\\
    & $c_j$ & $0$ & $0$ & $0$ & $0$ & $0$ & $0$ & $-1$ & \\
    \midrule
    $0$ & $y_1$ & $1$ & $-\tfrac{3}{2}$ & $0$ & $0$ & $-\tfrac{1}{4}$ & $-\tfrac{3}{2}$ & $1$ & $\tfrac{31}{4}$\\
    $0$ & $y_3$ & $0$ & $-2$ & $1$ & $0$ & $-\tfrac{1}{2}$ & $-1$ & $1$ & $\tfrac{15}{2}$\\
    $0$ & $y_4$ & $0$ & $-1$ & $0$ & $1$ & $0$ & $-1$ & $1$ & $9$\\
    \midrule
    & $z_j$ & $0$ & $0$ & $0$ & $0$ & $0$ & $0$ & $0$ & $0$\\
    & $z_j-c_j$ & $0$ & $0$ & $0$ & $0$ & $0$ & $0$ & $1$ & $0$\\
    \bottomrule
  \end{tabular}
\end{center}

The Phase I value is zero, with $a_2$ nonbasic and zero. Delete its column;
retain the nonsingular basis $(y_1,y_3,y_4)$.

\subsection{Phase II and absence of alternative optima}
Recompute the objective from $(-2,1,-5,1,0,0)$; do not reuse the Phase I row.

\begin{center}
  \small
  \renewcommand{\arraystretch}{1.25}
  \begin{tabular}{r@{\hspace{0.2em}}|@{\hspace{0.2em}}c@{\enspace}|@{\enspace}rrrrrr@{\quad}|@{\quad}r}\toprule
    $c_B$ & Basic & $y_1$ & $y_2$ & $y_3$ & $y_4$ & $s_1$ & $s_3$ & RHS\\
    & $c_j$ & $-2$ & $1$ & $-5$ & $1$ & $0$ & $0$ & \\
    \midrule
    $-2$ & $y_1$ & $1$ & $-\tfrac{3}{2}$ & $0$ & $0$ & $-\tfrac{1}{4}$ & $-\tfrac{3}{2}$ & $\tfrac{31}{4}$\\
    $-5$ & $y_3$ & $0$ & $-2$ & $1$ & $0$ & $-\tfrac{1}{2}$ & $-1$ & $\tfrac{15}{2}$\\
    $1$ & $y_4$ & $0$ & $-1$ & $0$ & $1$ & $0$ & $-1$ & $9$\\
    \midrule
    & $z_j$ & $-2$ & $12$ & $-5$ & $1$ & $3$ & $7$ & $-44$\\
    & $z_j-c_j$ & $0$ & $11$ & $0$ & $0$ & $3$ & $7$ & $-44$\\
    \bottomrule
  \end{tabular}
\end{center}

All indicators are nonnegative, so no Phase II pivot is needed:
\[
  y^*=(31/4,0,15/2,9),\qquad b^Ty^*=-44.
\]
For every feasible dual, the tableau implies
\begin{equation}\label{eq:task5b-bound}
  b^Ty=-44-11y_2-3s_1-7s_3\leq-44.
\end{equation}
A distinct optimizer would have to make an actual feasible move keeping the
objective unchanged. Equality forces $y_2=s_1=s_3=0$; the three basic equations
then fix $(y_1,y_3,y_4)=(31/4,15/2,9)$. Thus there is no such move and the dual
optimum is unique. The zero indicators belong only to basic variables; they
do not alone signal alternative optima.

\subsection{Recover the primal without a primal simplex solve}
Since $y_1,y_3,y_4>0$, complementary slackness makes primal rows 1, 3, 4 tight:
\[
  -2x_2-2x_3=-2,\qquad -2x_1+x_2+x_3=-5,\qquad
  2x_1+2x_2+x_3=1.
\]
The first gives $x_2+x_3=1$. Substitution into the second gives $x_1=3$;
the third becomes $2x_2+x_3=-5$. Therefore
\[
  x^*=(3,-6,7),\quad Ax^*-b=(0,11,0,0),\quad
  c^Tx^*=9-60+7=-44.
\]
The free variable $x_2=-6$ is admissible. Both constrained primal variables are
nonnegative; the unused row 2 has surplus 11. All dual column slacks are zero.
Feasibility and equal objectives certify both optima. The three equations above
are independent, so the primal optimum is unique as well.

\section{Task 5(c): equality domains reveal infeasibility}\label{sec:task5c}
Both $x_1$ and $x_2$ are free, giving two dual equalities:
\[
  \begin{aligned}
    \max\quad &-y_1+5y_2+y_3-y_4,\\
    -y_1-y_2-y_3&=-1,\\
    -y_1+y_3-y_4&=5,\\
    2y_1+3y_2+y_3-2y_4&\leq8,\qquad y\geq0.
  \end{aligned}
\]
The first equality implies $y_1+y_2+y_3=1$, hence
$-y_1+y_3-y_4\leq y_3\leq1$. It cannot equal 5.

\subsection{One Phase I pivot makes the contradiction quantitative}
Normalize the first equality by $-1$, append artificials $a_1,a_2$ to the
first two rows, and use slack $s_3$ in the third. Maximize $-a_1-a_2$ with
$B=(a_1,a_2,s_3)$.

\begin{center}
  \small
  \renewcommand{\arraystretch}{1.25}
  \begin{tabular}{r@{\hspace{0.2em}}|@{\hspace{0.2em}}c@{\enspace}|@{\enspace}rrrrrrr@{\quad}|@{\quad}r@{\qquad\quad}l@{\quad}c}\cmidrule(r{1.5em}){1-10}
    $c_B$ & Basic & $y_1$ & $y_2$ & $y_3$ & $y_4$ & $s_3$ & $a_1$ & $a_2$ & RHS & Ratio ($y_3\uparrow$) & Out\\
    & $c_j$ & $0$ & $0$ & $0$ & $0$ & $0$ & $-1$ & $-1$ & & & \\
    \cmidrule(r{1.5em}){1-10}
    $-1$ & $a_1$ & $1$ & $1$ & $1$ & $0$ & $0$ & $1$ & $0$ & $1$ & $1 / 1=1$ & $\longrightarrow$\\
    $-1$ & $a_2$ & $-1$ & $0$ & $1$ & $-1$ & $0$ & $0$ & $1$ & $5$ & $5 / 1=5$ & \\
    $0$ & $s_3$ & $2$ & $3$ & $1$ & $-2$ & $1$ & $0$ & $0$ & $8$ & $8 / 1=8$ & \\
    \cmidrule(r{1.5em}){1-10}
    & $z_j$ & $0$ & $-1$ & $-2$ & $1$ & $0$ & $-1$ & $-1$ & $-6$ & & \\
    & $z_j-c_j$ & $0$ & $-1$ & $-2$ & $1$ & $0$ & $0$ & $0$ & $-6$ & & \\
    \cmidrule(r{1.5em}){1-10}
  \end{tabular}
\end{center}

Enter $y_3$, whose indicator is $-2$, and remove $a_1$ after step 1.

\begin{center}
  \small
  \renewcommand{\arraystretch}{1.25}
  \begin{tabular}{r@{\hspace{0.2em}}|@{\hspace{0.2em}}c@{\enspace}|@{\enspace}rrrrrrr@{\quad}|@{\quad}r}\toprule
    $c_B$ & Basic & $y_1$ & $y_2$ & $y_3$ & $y_4$ & $s_3$ & $a_1$ & $a_2$ & RHS\\
    & $c_j$ & $0$ & $0$ & $0$ & $0$ & $0$ & $-1$ & $-1$ & \\
    \midrule
    $0$ & $y_3$ & $1$ & $1$ & $1$ & $0$ & $0$ & $1$ & $0$ & $1$\\
    $-1$ & $a_2$ & $-2$ & $-1$ & $0$ & $-1$ & $0$ & $-1$ & $1$ & $4$\\
    $0$ & $s_3$ & $1$ & $2$ & $0$ & $-2$ & $1$ & $-1$ & $0$ & $7$\\
    \midrule
    & $z_j$ & $2$ & $1$ & $0$ & $1$ & $0$ & $1$ & $-1$ & $-4$\\
    & $z_j-c_j$ & $2$ & $1$ & $0$ & $1$ & $0$ & $2$ & $0$ & $-4$\\
    \bottomrule
  \end{tabular}
\end{center}

All indicators are nonnegative: Phase I stops at $z_{\mathrm I}=-4$, hence
$w=a_1+a_2=4$. Indeed its second equation reads
\[
  -2y_1-y_2-y_4-a_1+a_2=4,
\]
so
\[
  a_1+a_2=4+2y_1+y_2+y_4+2a_1\geq4.
\]
This supplies a hand certificate of the Phase I optimum and the original
dual's infeasibility. There is no valid Phase II or recovered optimal primal.

The primal nevertheless has the feasible ray
\[
  x(t)=(-5-t,-t,0),\qquad t\geq0.
\]
Its row values are $(5+2t,5+t,5,t)$, dominating $(-1,5,1,-1)$,
and its objective is $5-4t\to-\infty$. Therefore the primal is unbounded below;
there are no dual optimal alternatives to discuss.

\section{Task 5(d): a feasible dual, two Phase II pivots}\label{sec:task5d}
Here only $x_2$ is free. The dual is
\[
  \begin{aligned}
    \max\quad &y_1+5y_2+y_3-2y_4,\\
    3y_1-y_2-y_3+3y_4&\leq1,\\
    y_1-y_3-y_4&=-5,\\
    2y_1+y_2+y_3-2y_4&\leq8,\qquad y\geq0.
  \end{aligned}
\]
Normalize the equality to $-y_1+y_3+y_4=5$ before adding artificial $a_2$;
use slacks $s_1,s_3$. The objective remains the original dual maximum.

\subsection{Phase I}
Start at $B=(s_1,a_2,s_3)$ and maximize $-a_2$.

\begin{center}
  \small
  \renewcommand{\arraystretch}{1.25}
  \begin{tabular}{r@{\hspace{0.2em}}|@{\hspace{0.2em}}c@{\enspace}|@{\enspace}rrrrrrr@{\quad}|@{\quad}r@{\qquad\quad}l@{\quad}c}\cmidrule(r{1.5em}){1-10}
    $c_B$ & Basic & $y_1$ & $y_2$ & $y_3$ & $y_4$ & $s_1$ & $s_3$ & $a_2$ & RHS & Ratio ($y_3\uparrow$) & Out\\
    & $c_j$ & $0$ & $0$ & $0$ & $0$ & $0$ & $0$ & $-1$ & & & \\
    \cmidrule(r{1.5em}){1-10}
    $0$ & $s_1$ & $3$ & $-1$ & $-1$ & $3$ & $1$ & $0$ & $0$ & $1$ & $\text{--}$ & \\
    $-1$ & $a_2$ & $-1$ & $0$ & $1$ & $1$ & $0$ & $0$ & $1$ & $5$ & $5 / 1=5$ & $\longrightarrow$\\
    $0$ & $s_3$ & $2$ & $1$ & $1$ & $-2$ & $0$ & $1$ & $0$ & $8$ & $8 / 1=8$ & \\
    \cmidrule(r{1.5em}){1-10}
    & $z_j$ & $1$ & $0$ & $-1$ & $-1$ & $0$ & $0$ & $-1$ & $-5$ & & \\
    & $z_j-c_j$ & $1$ & $0$ & $-1$ & $-1$ & $0$ & $0$ & $0$ & $-5$ & & \\
    \cmidrule(r{1.5em}){1-10}
  \end{tabular}
\end{center}

Indicators $q_3=q_4=-1$ tie; take $y_3$ by its smaller index.
Remove $a_2$ after step 5.

\begin{center}
  \small
  \renewcommand{\arraystretch}{1.25}
  \begin{tabular}{r@{\hspace{0.2em}}|@{\hspace{0.2em}}c@{\enspace}|@{\enspace}rrrrrrr@{\quad}|@{\quad}r}\toprule
    $c_B$ & Basic & $y_1$ & $y_2$ & $y_3$ & $y_4$ & $s_1$ & $s_3$ & $a_2$ & RHS\\
    & $c_j$ & $0$ & $0$ & $0$ & $0$ & $0$ & $0$ & $-1$ & \\
    \midrule
    $0$ & $s_1$ & $2$ & $-1$ & $0$ & $4$ & $1$ & $0$ & $1$ & $6$\\
    $0$ & $y_3$ & $-1$ & $0$ & $1$ & $1$ & $0$ & $0$ & $1$ & $5$\\
    $0$ & $s_3$ & $3$ & $1$ & $0$ & $-3$ & $0$ & $1$ & $-1$ & $3$\\
    \midrule
    & $z_j$ & $0$ & $0$ & $0$ & $0$ & $0$ & $0$ & $0$ & $0$\\
    & $z_j-c_j$ & $0$ & $0$ & $0$ & $0$ & $0$ & $0$ & $1$ & $0$\\
    \bottomrule
  \end{tabular}
\end{center}

The Phase I value is zero. Delete the nonbasic artificial.
The restored objective coefficients are $(1,5,1,-2,0,0)$.

\subsection{Phase II}

\begin{center}
  \small
  \renewcommand{\arraystretch}{1.25}
  \begin{tabular}{r@{\hspace{0.2em}}|@{\hspace{0.2em}}c@{\enspace}|@{\enspace}rrrrrr@{\quad}|@{\quad}r@{\qquad\quad}l@{\quad}c}\cmidrule(r{1.5em}){1-9}
    $c_B$ & Basic & $y_1$ & $y_2$ & $y_3$ & $y_4$ & $s_1$ & $s_3$ & RHS & Ratio ($y_2\uparrow$) & Out\\
    & $c_j$ & $1$ & $5$ & $1$ & $-2$ & $0$ & $0$ & & & \\
    \cmidrule(r{1.5em}){1-9}
    $0$ & $s_1$ & $2$ & $-1$ & $0$ & $4$ & $1$ & $0$ & $6$ & $\text{--}$ & \\
    $1$ & $y_3$ & $-1$ & $0$ & $1$ & $1$ & $0$ & $0$ & $5$ & $\text{--}$ & \\
    $0$ & $s_3$ & $3$ & $1$ & $0$ & $-3$ & $0$ & $1$ & $3$ & $3 / 1=3$ & $\longrightarrow$\\
    \cmidrule(r{1.5em}){1-9}
    & $z_j$ & $-1$ & $0$ & $1$ & $1$ & $0$ & $0$ & $5$ & & \\
    & $z_j-c_j$ & $-2$ & $-5$ & $0$ & $3$ & $0$ & $0$ & $5$ & & \\
    \cmidrule(r{1.5em}){1-9}
  \end{tabular}
\end{center}

Enter $y_2$, with indicator $-5$, and remove $s_3$ after step 3.
The objective rises from 5 to 20.

\begin{center}
  \small
  \renewcommand{\arraystretch}{1.25}
  \begin{tabular}{r@{\hspace{0.2em}}|@{\hspace{0.2em}}c@{\enspace}|@{\enspace}rrrrrr@{\quad}|@{\quad}r@{\qquad\quad}l@{\quad}c}\cmidrule(r{1.5em}){1-9}
    $c_B$ & Basic & $y_1$ & $y_2$ & $y_3$ & $y_4$ & $s_1$ & $s_3$ & RHS & Ratio ($y_4\uparrow$) & Out\\
    & $c_j$ & $1$ & $5$ & $1$ & $-2$ & $0$ & $0$ & & & \\
    \cmidrule(r{1.5em}){1-9}
    $0$ & $s_1$ & $5$ & $0$ & $0$ & $1$ & $1$ & $1$ & $9$ & $9 / 1=9$ & \\
    $1$ & $y_3$ & $-1$ & $0$ & $1$ & $1$ & $0$ & $0$ & $5$ & $5 / 1=5$ & $\longrightarrow$\\
    $5$ & $y_2$ & $3$ & $1$ & $0$ & $-3$ & $0$ & $1$ & $3$ & $\text{--}$ & \\
    \cmidrule(r{1.5em}){1-9}
    & $z_j$ & $14$ & $5$ & $1$ & $-14$ & $0$ & $5$ & $20$ & & \\
    & $z_j-c_j$ & $13$ & $0$ & $0$ & $-12$ & $0$ & $5$ & $20$ & & \\
    \cmidrule(r{1.5em}){1-9}
  \end{tabular}
\end{center}

Now enter $y_4$, with indicator $-12$, and remove $y_3$ after step 5. The objective rises by $12\cdot5=60$.

\begin{center}
  \small
  \renewcommand{\arraystretch}{1.25}
  \begin{tabular}{r@{\hspace{0.2em}}|@{\hspace{0.2em}}c@{\enspace}|@{\enspace}rrrrrr@{\quad}|@{\quad}r}\toprule
    $c_B$ & Basic & $y_1$ & $y_2$ & $y_3$ & $y_4$ & $s_1$ & $s_3$ & RHS\\
    & $c_j$ & $1$ & $5$ & $1$ & $-2$ & $0$ & $0$ & \\
    \midrule
    $0$ & $s_1$ & $6$ & $0$ & $-1$ & $0$ & $1$ & $1$ & $4$\\
    $-2$ & $y_4$ & $-1$ & $0$ & $1$ & $1$ & $0$ & $0$ & $5$\\
    $5$ & $y_2$ & $0$ & $1$ & $3$ & $0$ & $0$ & $1$ & $18$\\
    \midrule
    & $z_j$ & $2$ & $5$ & $13$ & $-2$ & $0$ & $5$ & $80$\\
    & $z_j-c_j$ & $1$ & $0$ & $12$ & $0$ & $0$ & $5$ & $80$\\
    \bottomrule
  \end{tabular}
\end{center}

This yields
\[
  y^*=(0,18,0,5),\qquad b^Ty^*=90-10=80.
\]
The final identity
\begin{equation}\label{eq:task5d-bound}
  b^Ty=80-y_1-12y_3-5s_3\leq80
\end{equation}
forces $y_1=y_3=s_3=0$ at every optimum. The remaining equations give
$y_4=5$, $y_2=18$, $s_1=4$. No nonzero feasible displacement can keep the
objective at 80, so the dual optimum is unique.

\subsection{Primal recovery and the role of the free variable}
The dual column slacks are
\[
  c-A^Ty^*=(4,0,0).
\]
Because $x_1\geq0$, complementarity gives $x_1=0$. Positive $y_2,y_4$ make
primal rows 2 and 4 tight:
\[
  -x_1+x_3=5,\qquad3x_1-x_2-2x_3=-2.
\]
Hence $x_3=5$, $x_2=-8$. The recovered point is
\[
  x^*=(0,-8,5),\qquad Ax^*-b=(1,0,12,0),\qquad
  c^Tx^*=40+40=80.
\]
It obeys the domains, both inactive inequalities have nonnegative surplus,
and the primal and dual values agree. These equations force the same point
for every primal optimum, proving its uniqueness without a primal simplex run.

\section{Interpretation and common mistakes}\label{sec:interpretation}
\begin{itemize}[leftmargin=*]
  \item \textbf{Objective signs:} task 4 negates the primal objective for its
        dual-simplex tableau; report 7 as the original minimum, not $-7$.
        Tasks 5 maximize the dual directly, including the valid negative maximum
        $-44$ in 5(b).
  \item \textbf{Free and nonpositive domains:} a free primal coordinate creates
        an equality in the dual. A nonpositive coordinate reverses the column
        inequality. Bounds on the right side do not determine these signs.
  \item \textbf{Two-phase transition:} artificials are feasibility devices.
        Delete their columns at a zero Phase I optimum and recompute the original
        objective. A negative Phase I maximum forbids Phase II.
  \item \textbf{Alternative optima:} a zero indicator is useful only with an
        actual feasible move. In 5(b) and 5(d), every nonbasic indicator is
        strictly positive, and the basic equations leave no optimal movement.
  \item \textbf{Primal recovery:} solve the equations forced by a dual optimizer,
        then verify all remaining inequalities and domains. Complementarity
        alone does not establish feasibility.
  \item \textbf{Unboundedness:} give a feasible starting point and an improving
        ray. An infeasible dual alone would not prove primal unboundedness if
        primal feasibility had not also been established.
\end{itemize}

\section{MATLAB comparison}\label{sec:matlab}
The companion
\href{run:matlab/exercises/sheet_01_dual_simplex_and_recovery.m}{\texttt{sheet\_01\_dual\_simplex\_and\_recovery.m}}
uses the same printed matrices, augmented equations, basis orders, and pivot
rules as the hand solutions. From the repository root, run
\begin{lstlisting}
run('matlab/exercises/sheet_01_dual_simplex_and_recovery.m')
\end{lstlisting}
It uses base MATLAB linear solves, with absolute tolerance $10^{-10}$ for
residuals, tableau comparisons, sign tests, and objective values. This tolerance
is appropriate to these small inputs; it is not a conditioning guarantee for
arbitrary models. The script checks the specified paths, including objective-row
recomputation, independently against new basis solves. It checks primal/dual
feasibility, complementarity, positive nonbasic indicators for uniqueness, and
feasible rays for the unbounded cases. It does not implement a general LP solver.

\begin{center}
  \begin{tabular}{llll}\toprule
    Task & Primal outcome & Dual outcome & Phase I artificial sum\\\midrule
    4 & Unique minimum $7$ & Unique maximum $7$ & Not used\\
    5(a) & Unbounded below & Infeasible & $3$\\
    5(b) & Unique minimum $-44$ & Unique maximum $-44$ & $0$\\
    5(c) & Unbounded below & Infeasible & $4$\\
    5(d) & Unique minimum $80$ & Unique maximum $80$ & $0$\\\bottomrule
  \end{tabular}
\end{center}
The table records exact expected hand results, not invented solver output.
Mathematical, static-analysis, and MATLAB execution evidence is identified
separately in the final remarks and repository verification record.

\section{Sources and editorial remarks}\label{sec:sources}
\noindent\textbf{Task source:} \texttt{../execrices-2026-27/ProgMatc1.pdf},
PDF page 2, tasks 4 and 5(a)--5(d). The Polish instructions and mathematical data
are preserved in Section~\ref{sec:statements}; layout is retypeset.
The supplied current-course context is 2026--27. The PDF creation metadata is
2025, which does not establish the delivery date. No lecture date is claimed.

The hand derivations, complete tableaus, recovery equations, rays, and uniqueness
arguments are explanatory additions. The generic instruction in task 5 requests
optimal solutions. Sections~\ref{sec:task5a} and~\ref{sec:task5c} show that the
printed domains and coefficients instead imply unbounded primal problems and
infeasible duals. The task statements remain unchanged; their intended variants
are unknown. Successful Phase II solves are provided for 5(b) and 5(d).

\textbf{Verification status:} hand tableaus and pivot choices were independently
recomputed using exact rational arithmetic. The user supplied successful MATLAB
output on 2026-10-11 and confirmed version 26.2.0.3386108 (R2026b).
All dual-simplex and recovery assertions passed, including the Phase I failures
and primal rays in 5(a) and 5(c). The supplied execution evidence is retained
in \path{docs/matlab-verification.md}.
\end{document}
